\documentclass[11pt]{article}
\usepackage[margin=1in]{geometry}
\usepackage{amsmath}
\usepackage{amssymb}
\usepackage{hyperref}
\usepackage{booktabs}
\usepackage{url}

\title{Memo 08: FVCOM-MP Effective Floc Size Coupling\\
       Using the ORIG\_SED Floc Settling Treatment}
\author{FVCOM-MP WaterPACT Investigation}
\date{May 2026}

\begin{document}
\sloppy
\maketitle

\section*{Purpose}

This memo records the technical implementation of a one-way FVCOM-MP coupling
to the original FVCOM sediment flocculation treatment.  The sediment module
remains unchanged.  The plastic module reconstructs the same local floc settling
signal used by \texttt{mod\_sed.F}, converts that settling speed to a
Stokes-equivalent floc diameter, and uses the result only inside the MP-floc
interaction pathway.

\section*{Theory Used}

\subsection*{ORIG\_SED floc settling speed}

For cohesive sediment, \texttt{mod\_sed.F} computes a local floc-enhanced
settling velocity in \texttt{Calc\_Wset\_Floc}.  The implemented empirical form
is
\begin{equation}
  W_s = W_0 F,
\end{equation}
where
\begin{equation}
  F =
  \frac{0.274 C^{0.48} S^{0.03} G^{0.22}}
       {(D_{50,\mathrm{mm}})^{0.58}} .
\end{equation}
Here $C$ is suspended sediment concentration in kg m$^{-3}$, $S$ is salinity,
$G$ is the turbulence/shear intensity argument described below, and
$D_{50,\mathrm{mm}}$ is the primary cohesive sediment diameter in mm.  The
existing ORIG\_SED guards are preserved:
\begin{itemize}
  \item if $F < 1$, set $F = 1$;
  \item if $S < 5$ or $S > 15$, set $F = 1$;
  \item cap the final settling speed at $W_s \le 0.5 \times 10^{-3}$ m s$^{-1}$.
\end{itemize}

\subsection*{G computation}

The MP implementation reconstructs the same $G$ expression used in
\texttt{mod\_sed.F}.  Element velocities are first mapped to nodes with
\texttt{E2N3D(u,unode)} and \texttt{E2N3D(v,vnode)}.  For each node and sigma
layer,
\begin{equation}
  j_e = \frac{m_f^2 (u_n^2 + v_n^2)}
             {\max(d,5)^{4/3}},
\end{equation}
with $m_f = 0.012$ and water depth $d$ in m.  The floc formula argument is then
\begin{equation}
  G =
  \left(
  \frac{g \sqrt{u_n^2+v_n^2} j_e}{V_{\mathrm{cons}}}
  \right)^{1/2},
\end{equation}
with $V_{\mathrm{cons}} = 10^{-6}$ and gravitational acceleration from the
FVCOM control module.  This is the same local calculation currently embedded in
the sediment deposition routine, including the 5 m shallow-depth floor.

\subsection*{Stokes-equivalent floc size}

The sediment model does not prognose floc diameter.  The new MP option therefore
uses an equivalent diameter obtained by assuming a spherical floc of constant
effective density:
\begin{equation}
  d_{\mathrm{floc,eff}} =
  \left[
  \frac{18 \mu W_s}
       {g(\rho_{\mathrm{floc,eff}}-\rho_w)}
  \right]^{1/2}.
\end{equation}
The parameter \texttt{PLAST\_FLOC\_RHO\_EFF} supplies
$\rho_{\mathrm{floc,eff}}$.  The default value is 1300 kg m$^{-3}$, matching the
legacy MP aggregation kernel's floc density.  The water dynamic viscosity is
computed from the same bounded temperature-dependent kinematic viscosity already
used in the MP aggregation kernel:
\begin{equation}
  \nu = 10^{-6} \exp[-0.025(T-20)],
  \qquad 0.5\times10^{-6} \le \nu \le 2.0\times10^{-6},
\end{equation}
and $\mu = \rho_w \nu$.  The implementation prevents the equivalent diameter
from falling below the primary sediment $D_{50}$.

\section*{Code Change Summary}

\begin{itemize}
  \item Added optional plastic inputs \texttt{PLAST\_FLOC\_EFFECTIVE} and
        \texttt{PLAST\_FLOC\_RHO\_EFF}.  Missing keys preserve the old behavior:
        the effective floc-size path is off and the density defaults to
        1300 kg m$^{-3}$.
  \item Added validation that \texttt{PLAST\_FLOC\_EFFECTIVE=T} requires
        \texttt{PLAST\_FLOC=T} and \texttt{-DORIG\_SED}.
  \item Added \texttt{Calc\_ORIG\_Floc\_Effective\_Cell} in
        \texttt{mod\_plast.F} to reconstruct ORIG\_SED $W_s$, compute $G$, and
        convert $W_s$ to $d_{\mathrm{floc,eff}}$.
  \item Updated MP aggregation so size exclusion, collision radius, sediment
        particle mass, and differential-settling kernels can use the effective
        floc diameter and reconstructed floc settling speed.
  \item Updated MP disaggregation so the maximum active floc size can come from
        the effective floc diameter field when the new switch is active.
  \item Updated aggregated MP settling so \texttt{wset\_floc} stores a
        sediment-concentration-weighted reconstructed floc settling speed.  The
        existing NetCDF variable \texttt{settle\_vel\_floc\_*} can therefore be
        used in post-processing to back out
        $d_{\mathrm{floc,eff}}$ with the same Stokes formula.
  \item Created \texttt{FVCOM\_MP\_test\_run/INPUT/generic\_plastic\_d.inp} and
        \texttt{FVCOM\_MP\_test\_run/INPUT/generic\_sediment\_d.inp} as the
        switch-on case.  The \texttt{\_c} files remain the comparison case.
\end{itemize}

\section*{Implementation Plan}

\begin{enumerate}
  \item Push the updated \texttt{mod\_plast.F} to the HPC source tree and compile
        with the existing FVCOM-MP build options plus \texttt{-DSEDIMENT},
        \texttt{-DORIG\_SED}, and \texttt{-DPLASTIC}.  The new MP option has a
        runtime guard, so a switch-on input will stop early if \texttt{ORIG\_SED}
        is absent.
  \item Keep the \texttt{\_c} inputs as the comparison case:
        \texttt{generic\_plastic\_c.inp} and
        \texttt{generic\_sediment\_c.inp}.  This case keeps
        \texttt{PLAST\_FLOC=T} but does not activate the effective floc-size
        option.
  \item Use the \texttt{\_d} inputs as the switch-on case:
        \texttt{generic\_plastic\_d.inp} sets
        \texttt{PLAST\_FLOC\_EFFECTIVE=T} and
        \texttt{PLAST\_FLOC\_RHO\_EFF=1300.0}; the new
        \texttt{generic\_sediment\_d.inp} switches the single sediment class to
        \texttt{SED\_TYPE=cohesive}, \texttt{SED\_SD50=0.004} mm, and
        \texttt{SED\_WSET=0.1} mm s$^{-1}$.  The sediment name is kept as
        \texttt{coarse\_sand\_1} so existing river-forcing variable names remain
        compatible.
  \item Create or copy a \texttt{RUN\_d} namelist from \texttt{RUN\_c} and point
        it to \texttt{generic\_sediment\_d.inp} and
        \texttt{generic\_plastic\_d.inp}.  Keep the hydrodynamic, river, source,
        and reporting settings aligned with \texttt{RUN\_c} unless a deliberate
        sensitivity change is being made.
  \item After a successful compile, start with a short HPC run of case
        \texttt{\_d}.  Confirm that the input guards pass, the run reaches the
        first few reports, and the existing \texttt{settle\_vel\_floc\_*}
        output is present for post-processing.
  \item Use the \texttt{\_c} and \texttt{\_d} outputs as a paired comparison.
        The \texttt{\_d} run is the ORIG\_SED effective-floc coupling case, and
        the \texttt{\_c} run remains the non-effective MP-floc comparison.
\end{enumerate}

\section*{Post-processing Note}

When \texttt{PLAST\_FLOC\_EFFECTIVE=T}, the output variable
\texttt{settle\_vel\_floc\_<plastic name>} stores the effective aggregated MP
floc-settling speed used by the plastic module.  A post-processing script can
recover the model-used equivalent floc diameter by applying the same Stokes
formula above with the output \texttt{ambient\_rho}, the configured
\texttt{PLAST\_FLOC\_RHO\_EFF}, and the same temperature-dependent viscosity
closure.  For multiple active sediment classes, this output represents the
sediment-concentration-weighted effective settling speed used for aggregated MP
settling, not a separate diameter for each sediment class.

\end{document}
